\chapter{宇宙論の歴史と最前線 --- ビッグバンからマルチバース、そしてAIが描く宇宙の地図}

\section*{本章の学習目標}
\begin{itemize}
\item アインシュタインの静的宇宙モデルから、ハッブルの「膨張する宇宙」へのパラダイムシフトを学ぶ。
\item ビッグバン理論の確立と、宇宙マイクロ波背景放射（CMB）の発見のドラマを知る。
\item ビッグバンの「前」に何が起きたか？ インフレーション理論が解決した宇宙の謎を理解する。
\item 宇宙の95\%を占める謎の存在「ダークマター」と「ダークエネルギー」の役割を知る。
\item 重力波の直接観測と、微小なノイズから宇宙の「音」を拾い上げるAIの活躍を学ぶ。
\item AIを用いて宇宙の大規模構造をスーパーコンピュータ内に再現する「デジタルツイン宇宙」の最前線を知る。
\item 系外惑星の発見からSETI（地球外知的生命体探査）に至る、AIが導くファーストコンタクトへの挑戦を学ぶ。
\item 「ホログラフィック原理」が突きつける、宇宙の本質は「情報」であるというパラダイムシフトを理解する。
\item マルチバース（多元宇宙）理論と、人間原理が突きつける哲学的な問いについて考える。
\end{itemize}

\section{静的な宇宙から「膨張する宇宙」へ：人類の宇宙観を覆すパラダイムシフト}

私たちが夜空を見上げるとき、星々は静かに、そして永遠にそこにあるように感じられます。古代ギリシャから20世紀の初めまで、人類の最も偉大な頭脳たちでさえ、「宇宙は無限の過去から無限の未来まで、大きさも姿も変わらない静的で不変なものである」と信じて疑いませんでした。

\subsection{永遠の宇宙への信仰と「オルバースのパラドックス」}

アイザック・ニュートンもまた、絶対時間と絶対空間の背景に、星々が静止して浮かぶ宇宙を想像していました。しかし、無限で静的な宇宙には、ある単純で致命的な矛盾が潜んでいました。それが19世紀に提起された\textbf{「オルバースのパラドックス」}です。

「もし宇宙が無限に広く、星が無限の過去から均等に散らばっているのなら、夜空のどの方向を見ても必ずどこかの星の表面にぶつかるはずだ。ならば、夜空は太陽の表面のように、全体がまばゆく光り輝いていなければならない。なぜ夜空は暗いのか？」

この素朴な疑問は、静的で無限な宇宙という常識が、実は根本的に間違っているかもしれないという最初の警告でした。

\subsection{アインシュタインの「最大の過ち」と宇宙項}

1915年、アルベルト・アインシュタインは重力の正体を「時空の歪み」として記述する\textbf{「一般相対性理論」}を完成させました。彼は自らの美しい方程式を、宇宙全体の形を計算するために適用しようと試みました。

しかし、計算結果は彼自身の信念を裏切るものでした。宇宙に散らばる星や銀河は、互いの重力（引力）で引き合っているため、放っておけばいつか一点に向かって潰れてしまう（収縮する）か、あるいは最初から猛烈な勢いで広がっているはずだ、という「動的な宇宙」の姿が方程式からあぶり出されたのです。

静的で永遠な宇宙を固く信じていたアインシュタインは、自らの方程式が示したこのダイナミズムを嫌悪しました。彼は宇宙がピタリと静止するように計算の辻褄を合わせるため、方程式の中に重力に逆らって空間を押し広げる謎の反発力\textbf{「宇宙項（宇宙定数：\(\Lambda\)）」}を人為的に、かつ無理やり付け足してしまったのです。

\subsection{異端の予言者たち：フリードマンとルメートル}

アインシュタインが自らの直感で数式を捻じ曲げたのに対し、「数式が示す真実」を純粋に信じた無名の天才たちがいました。

1922年、ロシアの数学者アレクサンドル・フリードマンは、宇宙項のないアインシュタイン方程式を素直に解き、「宇宙は膨張するか収縮するかのどちらかである」と数学的に証明しました。さらに1927年、ベルギーの司祭であり物理学者のジョルジュ・ルメートルも独立して「宇宙は膨張している」という理論を導き出し、遠くの銀河ほど速く遠ざかっているはずだと予言しました。

しかし、当時のアインシュタインは彼らの論文を読み、「君の数学は正しいが、物理学としてはおぞましい（abominable）」と酷評し、全く取り合おうとしませんでした。宇宙が動いているなどという不条理は、当時の物理学界の頂点に立つ彼にとってすら、受け入れがたいものだったのです。

\subsection{ハッブルの巨大望遠鏡と「赤方偏移」の衝撃}

この机上の空論の対立を、圧倒的な「観測事実（データ）」によって粉々に打ち砕いたのが、アメリカの天文学者エドウィン・ハッブルです。

1920年代、ハッブルは当時世界最大を誇るウィルソン山天文台のフッカー望遠鏡を使い、夜空の「星雲」と呼ばれていたモヤモヤした光のシミを観測しました。彼は、女性天文学者ヘンリエッタ・スワン・リービットが発見していた「セファイド変光星（明るさの変化の周期で真の明るさが分かる、宇宙の灯台）」を利用して星雲までの距離を測り、それらが私たちの天の川銀河の外にある「別の巨大な銀河（島宇宙）」であることを人類で初めて証明しました。

さらに彼は、それらの銀河から届く光の波長を分光器で分析しました。救急車が遠ざかるときにサイレンの音が低くなる（ドップラー効果）のと同じように、光を出している天体が地球から遠ざかっていると、光の波長が引き伸ばされて「赤っぽく」見えます。これを\textbf{「赤方偏移（Redshift）」}と呼びます。

ハッブルの観測結果は、世界の常識を根底から覆すものでした。多くの銀河が赤方偏移を示し、地球から遠ざかっていたのです。しかも、ルメートルが予言した通り、「遠くにある銀河ほど、より速いスピードで遠ざかっている（ハッブル＝ルメートルの法則）」ことが観測的に示されました。

\subsection{空間そのものが伸びている}

これは、地球が宇宙の中心にあって、皆が地球を嫌って逃げていくという意味ではありません。風船の表面にマジックで無数の銀河を描き、風船を膨らませると、どの銀河から見ても他のすべての銀河が遠ざかっていくように見えます。銀河が空間の中を移動しているのではなく、\textbf{「銀河と銀河の間にある『空間そのもの』が、今この瞬間も膨張し続けている」}という決定的な証拠だったのです。

この圧倒的な観測データを突きつけられたアインシュタインは、ついに自らの誤りを認めざるを得ませんでした。1931年、彼はハッブルの天文台を訪れて望遠鏡を覗き込み、静的な宇宙を捏造するために方程式に宇宙項を付け足したことを\textbf{「生涯最大の過ち（My biggest blunder）」}と呼んで深く後悔し、自らの手で宇宙項を削除しました。

こうして、永遠で不変だと思われていた宇宙は、ダイナミックに成長し続ける「膨張する宇宙」へと、人類史上最大のパラダイムシフトを果たしたのです。そしてこの事実は、直ちに次なる究極の問いを生み出すことになります。「宇宙が今も膨張し続けているのなら、時間を過去へと巻き戻したとき、一体何が起こるのか？」

\section{ビッグバン理論の勝利と「宇宙の産声」}

宇宙が今も膨張し続けているのなら、時間を過去へと巻き戻せば、宇宙はどんどん小さくなり、最終的には「想像を絶するほど超高温・超高密度の極小の一点」に収束するはずです。1931年、ルメートルはこの始まりの一点を「原始的原子（Primeval Atom）」と呼び、そこからの爆発的な膨張によって宇宙が始まったと主張しました。

\subsection{「ビッグバン」という皮肉な命名と定常宇宙論}

しかし、この「宇宙に始まり（誕生の瞬間）がある」という考え方は、当時の多くの物理学者から激しい拒絶に遭いました。「始まりの前には何があったのか」という宗教的・形而上学的な匂いが強すぎたためです。

対立する学者たちは、\textbf{「定常宇宙論」}を提唱しました。これは「宇宙は確かに膨張しているが、隙間ができた分だけ『無から新しい物質』が湧き出してくるため、宇宙全体の密度や姿は永遠に変わらない」という説です。

定常宇宙論の強力なリーダーであったイギリスの天体物理学者フレッド・ホイルは、1949年のBBCラジオ番組に出演した際、敵対するルメートルらの理論を馬鹿にしてこう言い放ちました。

「彼らの言うように、宇宙が過去のある一点の『大爆発（Big Bang）』から始まっただなんて、実に滑稽な話だ」

ところが歴史の皮肉なことに、ホイルが揶揄するために使ったこの\textbf{「ビッグバン」}というキャッチーな言葉が、大衆の心を掴み、そのまま理論の正式名称として定着してしまったのです。

\subsection{宇宙の料理鍋：最初の3分間と元素の起源}

このビッグバン理論を強力に推し進め、決定的な「物理的証拠」を突きつけたのが、ロシア出身の物理学者ジョージ・ガモフとその弟子のラルフ・アルファーでした（1948年）。

当時、宇宙を観測すると「水素が約75\%、ヘリウムが約25\%」という質量の比率で満たされていることが分かっていました。なぜこの比率なのか？ ガモフは、アインシュタインの相対性理論と最新の原子核物理学を使い、ビッグバン直後の超高温宇宙を「巨大な料理鍋」に見立てて計算を行いました。

宇宙誕生から約1秒後、宇宙は100億度という想像を絶する高温でした。そこから\textbf{「わずか3分間」の間に、宇宙が膨張して温度が10億度まで下がる過程で、陽子と中性子が激しく衝突し合い、次々と核融合を起こしました。ガモフらがこの「3分間の料理プロセス」を計算した結果、生成される物質の比率はおおよそ「水素75\%、ヘリウム25\%」となり、観測される宇宙の元素組成をよく説明できることが分かったのです。 この有名な論文は、ガモフの悪ふざけで友人のハンス・ベーテの名前を勝手に共著者に加え、ギリシャ文字の最初の3文字にちなんで「\(\alpha\beta\gamma\)（アルファー・ベーテ・ガモフ）理論」}と呼ばれ、ビッグバン理論の圧倒的な証拠となりました。

\subsection{光の脱獄：「宇宙の晴れ上がり」と残光の予言}

ガモフはさらに、もう一つのとんでもない予言をしました。

「もし宇宙が超高温の火の玉から始まったのなら、その当時の猛烈な光（熱放射）の残骸が、今でも宇宙を漂っているはずだ」というものです。

ビッグバン直後の宇宙は、電子と原子核がバラバラに飛び交う「プラズマ状態」でした。光（フォトン）は電子にぶつかって真っ直ぐ進むことができず、当時の宇宙はまるで「濃い霧」に包まれたように真っ暗でした。

しかし、ビッグバンから約38万年後、宇宙が膨張して温度が約3000度まで下がると、劇的な変化が起きます。飛び回っていた電子が原子核の引力に捕獲され、電気的に中性な「原子」が初めて誕生したのです。電子という障害物が空間から消え去った瞬間、それまで閉じ込められていた光が一斉に、無限の宇宙空間へと真っ直ぐに解き放たれました。これを\textbf{「宇宙の晴れ上がり（Recombination）」}と呼びます。

ガモフは予言しました。

「その38万年目に放たれた約3000度の光は、その後138億年にわたる宇宙の膨張によって波長が引き伸ばされ、現在では温度が劇的に下がり、約マイナス270度（数ケルビン）の極めて冷たいマイクロ波（電波）として、宇宙のあらゆる方向から等しく降り注いでいるはずだ」

しかし、当時の観測技術ではそのような微弱な電波を捉えることはできず、ガモフの偉大な予言は次第に忘れ去られていきました。

\subsection{鳩のフンと世紀のすれ違い：CMBの偶然の発見}

ガモフの予言から約20年後の1964年。アメリカのベル研究所で、通信用の巨大なホーン型アンテナを調整していた二人の研究者、アーノ・ペンジアスとロバート・ウィルソンは、奇妙な現象に悩まされていました。

アンテナを空のどの方向に向けても、季節を変えても、昼でも夜でも、常に「ザーッ」という謎の雑音（マイクロ波のバックグラウンドノイズ）が受信されるのです。彼らは軍のレーダーの干渉や、ニューヨークの都市ノイズ、さらにはアンテナに巣食った鳩のフンが原因だと思い、自ら鳩を追い払い懸命に掃除をしましたが、ノイズは決して消えませんでした。

困り果てた二人は、わずか数十キロ離れたプリンストン大学の研究者に相談の電話をかけました。

その電話を受けたロバート・ディッケらプリンストン大のチームは、受話器を置いた瞬間に「我々は出し抜かれた（We've been scooped!）」と叫んだと伝えられています。実はディッケのチームは、忘れ去られていたガモフの理論に自力でたどり着き、「ビッグバンの残光」を見つけるための専用のアンテナをまさに建設中だったのです。

ペンジアスとウィルソンが鳩のフンのせいだと思って悩まされていたこの全天から降り注ぐ約マイナス270度（3ケルビン）の電波こそが、\textbf{ガモフが予言したビッグバン開始38万年目の最初の光「宇宙マイクロ波背景放射（CMB：Cosmic Microwave Background）」}だったのです。

この「宇宙の産声」の偶然の発見により、定常宇宙論は大きく退潮し、ビッグバン理論は宇宙論の標準的な枠組みとなりました（ペンジアスとウィルソンはこの功績でノーベル物理学賞を受賞します）。

\section{インフレーション：究極の「無」からの創生とビッグバンの起爆剤}

CMBの発見によってビッグバン理論は確固たるものになりました。しかし、観測技術が向上し、宇宙の姿がより詳細に分かってくると、物理学者たちはビッグバン理論だけでは絶対に説明できない\textbf{「3つの致命的な謎」}に直面することになります。

\subsection{ビッグバン理論が抱えた3つの致命的な謎}

地平線問題：CMBの温度を精密に測ると、宇宙の「右の果て」と「左の果て」で、温度が10万分の1の精度で完全に一致（約2.725ケルビン）していることが分かりました。しかし、右の果てと左の果てはあまりにも遠く離れているため、宇宙が誕生してから現在までの間に、光の速さをもってしても情報（熱）を伝えることができません。一度も連絡を取り合っていないはずの遠く離れた場所同士が、なぜ示し合わせたかのように「全く同じ温度」になれたのでしょうか？

平坦性問題：アインシュタインの一般相対性理論によれば、宇宙の空間は物質の重力によって「丸く閉じる」か「鞍のように反り返る」可能性があります。しかし観測データは、私たちの宇宙空間が非常に平坦に近いことを示していました。なぜこれほどまでに平坦さが保たれているのでしょうか？

モノポール（磁気単極子）問題：第13章で触れた大統一理論（GUT）によれば、ビッグバン直後の超高温状態では、N極だけ、あるいはS極だけの性質を持つ「磁気単極子（モノポール）」という極めて重い素粒子が大量に生成されるはずです。しかし、現在までのところ宇宙のどこを探してもモノポールはただの1個も見つかっていません。大量に生まれたはずのモノポールはどこへ消えたのでしょうか？

\subsection{空間の指数関数的膨張：「インフレーション理論」の誕生}

この3つの致命的な謎を、たった一つのアイデアで一挙に解決する候補として登場したのが、1980年代初頭に日本の佐藤勝彦とアメリカのアラン・グースがそれぞれ独立に提唱した\textbf{「インフレーション理論」}です。

インフレーション理論は、ビッグバンの「火の玉」ができる直前、宇宙が誕生した最初の\(10^{-38}\)秒という想像を絶する一瞬の間に、空間そのものが\textbf{「光速をはるかに超えるスピードで、指数関数的に（数十桁も）急膨張した」}と予言しました（※物質が光速を超えることはできませんが、空間そのものが膨張するスピードに上限はないため、相対性理論には反しません）。

この急膨張の凄まじさは、「ウイルスほどの大きさだった宇宙が、一瞬で銀河系の大きさにまで引き伸ばされた」ようなスケールです。

この「一瞬の超急膨張」を仮定するだけで、すべての謎が美しく解けました。

地平線問題の解決：現在遠く離れている右の果てと左の果ては、インフレーションが起きる前は「極微のミクロな一つの領域」の中に密着しており、十分に熱をやり取りして同じ温度になっていました。それがインフレーションによって一瞬で宇宙の果てと果てに引き離されたため、現在観測しても同じ温度になっているのです。

平坦性問題の解決：丸い風船も、一気に地球サイズまで膨らませれば、表面にいるアリにとっては「ほとんど平ら」に見えます。宇宙空間も、想像を絶する急膨張によって極限まで引き伸ばされたため、観測できる範囲では非常に平坦に見えるようになったのです。

モノポール問題の解決：ビッグバン直後に生まれた大量のモノポールも、空間そのものが数十桁も急膨張したため、宇宙の彼方へ極限まで「薄められて」しまいました。そのため、現在の私たちの観測範囲の中には1個も見当たらないほど希少な存在になってしまったのです。

\subsection{「偽の真空」からの相転移と、ビッグバンの「着火」}

では、何がそのような凄まじい空間の急膨張を引き起こし、そしてなぜそれは終わったのでしょうか？ その原動力こそが、第9章で学んだ\textbf{「相転移（Phase Transition）」と、量子力学における「真空のエネルギー」}です。

インフレーション理論によれば、誕生したばかりの極微の宇宙は、何もないように見えて実は膨大なエネルギーをパンパンに抱え込んでいる\textbf{「偽の真空（False Vacuum）」}と呼ばれる不安定な状態にありました。この「偽の真空」が持つ猛烈な反発力（負の圧力）が、空間の指数関数的な急膨張（インフレーション）を引き起こしたのです。

しかし、過冷却状態の水に衝撃を与えると一瞬で氷になるように、この不安定な「偽の真空」は、ある瞬間にガクンとエネルギーの低い\textbf{「真の真空（True Vacuum）」}へと状態を変化させました（これを真空の相転移と呼びます）。

相転移が起きた瞬間、急膨張はピタリと止まります。そして、水が氷になるときに潜熱を放出するように、「偽の真空」が持っていた膨大なエネルギーが一気に解放され、冷え切っていた広大な宇宙空間をドカンと強烈に再加熱（リヒーティング）したのです。

この莫大な熱エネルギーによって、空間に無数の素粒子（クォークや電子など）が産み落とされたと考えられています。これが、ガモフが想像した\textbf{「超高温の火の玉（ビッグバン）」を生み出した過程だと解釈できます。 つまり、現代的な見方では「インフレーションによる急膨張\(\to\)真空の相転移\(\to\)再加熱\(\to\)熱いビッグバン」}という順番で宇宙初期を理解します。

\subsection{量子ゆらぎの凍結：星と銀河の「種」}

さらに驚くべきことに、インフレーション理論は「なぜ宇宙に星や銀河が存在するのか」という根本的な問いにも、有力な説明を与えました。

第7章で学んだ通り、ミクロな量子力学の世界では「不確定性原理」により、エネルギーが常にデタラメに波立ち、揺らいでいます（量子ゆらぎ）。インフレーションの超急膨張は、このミクロな\textbf{「ただの確率的な量子のノイズ（ゆらぎ）」を、一瞬にしてマクロな宇宙スケールにまで引き伸ばし、空間にわずかな「濃淡（密度の違い）」として凍りつかせた}のです。

その証拠は、まさにCMB（宇宙マイクロ波背景放射）の中に刻まれていました。1990年代以降、COBEやWMAP、プランクといった観測衛星がCMBを精密に観測したところ、ほぼ均一だと思われていた温度の中に「10万分の1度」というごくわずかなムラ（ゆらぎ）の模様が発見されました。

このわずかに温度が高い（密度が高い）部分は、周囲よりもほんの少しだけ重力が強いため、138億年という長い時間をかけて周囲のガスを引き寄せて星となり、やがて銀河、そして銀河団へと成長していきました。

私たち人間や地球、そして夜空に輝く無数の銀河の起源は、宇宙誕生直後のミクロな「量子のノイズ」が、インフレーションという巨大な虫眼鏡でマクロな世界に拡大コピーされたものだったのです。

\section{暗黒の支配者：ダークマターとダークエネルギーの衝撃}

インフレーション理論とビッグバン理論により、人類はついに宇宙の誕生から現在に至る壮大な進化のシナリオを手に入れました。しかし、1970年代から21世紀にかけての観測技術の飛躍的な進歩は、私たちをかつてない絶望の淵へと突き落とします。

私たちが知っている物理学（クォークや電子など）で作られた星やガスは、宇宙全体の\textbf{わずか「5\%」}に過ぎなかったのです。

\subsection{ツビッキーの孤独な警告と「かみのけ座銀河団」}

宇宙に「見えない重力源」が存在するかもしれないという最初の警告は、実は1933年に遡ります。スイスの天文学者フリッツ・ツビッキーは、数千個の銀河が群れをなす「かみのけ座銀河団」を観測し、銀河同士が猛烈なスピードで動き回っていることに気づきました。

彼が計算したところ、目に見えている星（銀河）の重力だけでは、これほど猛スピードで動く銀河団を繋ぎ止めておくことはできず、宇宙空間にバラバラに四散してしまうはずでした。彼は「光を出さない未知の物質（Dunkle Materie：ダークマター）が、見えている質量の400倍以上存在しなければならない」と主張しました。しかし、彼の主張はあまりに突飛すぎたため、当時の学界からは完全に無視されてしまいました。

\subsection{ヴェラ・ルービンと「銀河回転曲線」の謎}

ツビッキーの警告から約40年後の1970年代、アメリカの女性天文学者ヴェラ・ルービンが、より精密で決定的な証拠を突きつけます。彼女は、個々の「渦巻き銀河」の中で星がどのように回転しているかを精密に観測しました。

ニュートン力学（ケプラーの法則）に従えば、太陽系の惑星と同じように、銀河の中心から遠く離れた縁の星ほど、ゆっくりと回るはずです。

しかし観測結果は驚くべきものでした。銀河の縁にある星々も、中心付近の星と全く同じような猛烈なスピードで回転していたのです。この猛スピードの遠心力に耐えて星が宇宙に投げ出されないようにするためには、\textbf{銀河全体をすっぽりと包み込む、見えない巨大な質量の塊（ダークマター・ハロ）}が絶対に必要でした。

\subsection{決定的な証拠「弾丸銀河団」と重力レンズ}

その後、アインシュタインの一般相対性理論が予言した「重力レンズ効果（重力で空間が歪み、背後の星の光が曲がって見える現象）」を使って、見えないダークマターの分布を直接マッピングする技術が確立されました。

そして2006年、ダークマターの存在を決定づける歴史的な観測データが得られました。それが\textbf{「弾丸銀河団（Bullet Cluster）」です。 これは、2つの巨大な銀河団が猛スピードで正面衝突した直後の姿を捉えたものです。通常の物質（高温のガス）は、ぶつかり合って摩擦を起こし、衝突の中央地点にドロドロに留まってX線を放っていました。 しかし、重力レンズで「質量の中心」を測定してみると、なんと質量の大部分は中央のガスを「完全にすり抜けて」、両外側へポツンと移動していたのです。 「重力は持っているが、他の物質と一切ぶつかったり摩擦を起こしたりしない透明な物質」。この弾丸銀河団の観測により、ダークマターの実在はもはや疑いようのない事実となりました。ダークマターは宇宙の約27\%}を占めており、第13章で述べたように、AIを駆使して世界中の加速器や地下実験施設でその正体（未知の素粒子であるWIMPやアクシオンなど）の探索が続けられています。

\subsection{宇宙の加速膨張と「ダークエネルギー」}

ダークマターの存在だけでも衝撃的でしたが、1998年、宇宙論を根底からひっくり返すさらに恐ろしい大発見がありました。

サウル・パールムッター率いるチームと、ブライアン・シュミット、アダム・リース率いるライバルチームの2つが、それぞれ独立して「Ia型超新星」という明るさが常に一定の星（宇宙の標準光源）を使い、宇宙の膨張速度の歴史を精密に測定しました。

当時の天文学者たちは誰もが、ビッグバンで膨張を始めた宇宙は、宇宙全体にある星やダークマターの「引力」によってブレーキがかかり、やがて「減速していく」と信じて疑いませんでした。彼らはただ「どれくらい減速しているか」を測ろうとしていたのです。

しかし結果は真逆でした。宇宙の膨張スピードは減速するどころか、約50億年前から\textbf{「どんどん加速（アクセルを踏み込んでいる）」していたのです。 空間そのものを外側へ向かって押し広げる、未知の猛烈な反発力。これこそが「暗黒エネルギー（ダークエネルギー）」です。ダークエネルギーは宇宙の総エネルギーの約68\%}という圧倒的な割合を占めており、今もこの瞬間も空間を無限の彼方へと引き裂こうとしています。

\subsection{アインシュタインの復活と「物理学史上最悪の予言」}

宇宙を押し広げる反発力。それはまさに、アインシュタインがかつて静止宇宙を保つための方程式に無理やり付け足し、のちに「生涯最大の過ち」として捨て去ったあの\textbf{「宇宙項（宇宙定数\(\Lambda\)）」}そのものでした。アインシュタインの「最大の過ち」は、実は宇宙を支配する究極の真理として、半世紀以上の時を経て見事に復活を遂げたのです。

しかし、このダークエネルギーの正体は、現代物理学を絶望の淵に立たせています。

最も有力な候補は、量子力学が予言する「真空のエネルギー（何もない空間そのものが持つエネルギー）」です。しかし、場の量子論の方程式を使って真空のエネルギー密度を素朴に見積もると、1998年に超新星の観測から得られた実際のダークエネルギーの密度よりも、なんと\textbf{「\(10^{120}\)倍（0が120個つくほど大きい）」という、途方もないズレが生じてしまうのです。 理論と観測がこれほどまでに絶望的に合わない問題は他に類を見ず、これは「物理学史上最悪の予言（宇宙項問題）」}と呼ばれ、現在に至るまで誰一人として解決できていません。

通常物質5\%、ダークマター27\%、ダークエネルギー68\%。

私たちが数千年の歴史をかけて築き上げてきた、ニュートン、マクスウェル、アインシュタイン、そして量子力学という完璧な物理学の城は、広大な宇宙のほんの「5\%」の波打ち際で遊んでいたに過ぎませんでした。残りの95\%は、いまだ人類の理解を完全に拒む絶対的な暗黒に包まれています。

\section{重力波天文学とAI：宇宙の「音」を聴く}

暗黒の宇宙を解き明かすために、21世紀に入り人類は全く新しい「宇宙を観測する目（あるいは耳）」を手に入れました。それが\textbf{「重力波（Gravitational Wave）」}です。

\subsection{時空のさざ波}

アインシュタインは100年前に、質量を持った物体が激しく動くと、その周囲の「時空（空間と時間）」そのものがゴムシートのように波立ち、光の速さで宇宙空間を伝わっていくと予言していました。

しかし、その時空の歪みはあまりにも微小です。例えば、太陽の数十倍の重さを持つブラックホール同士が衝突したとしても、地球に届く頃には、\textbf{「地球から太陽までの距離が、水素原子1個分だけ伸び縮みする」}という、想像を絶するほど小さな歪みになってしまいます。アインシュタイン自身も「人類がこれを直接観測することは永遠に不可能だろう」と考えていました。

\subsection{LIGOの奇跡と機械学習の威力}

2015年、アメリカの巨大なレーザー干渉計「LIGO（ライゴ）」が、13億光年彼方で2つのブラックホールが合体した瞬間に発せられた重力波を、人類史上初めて直接捉えることに成功しました。

この奇跡的な観測の裏には、AIと機械学習の圧倒的な貢献があります。

LIGOの検出器には、地面の微振動、遠くを走るトラックの振動、空気のゆらぎなど、宇宙からの信号の何百万倍もの巨大な「ノイズ」が絶えず入り込んでいます。このノイズの嵐の中から、ブラックホールが合体する瞬間に発する特有の波形（チャープ信号：鳥のさえずりのように周波数が上がる波形）を見つけ出すのは至難の業です。

ここで、現代の天体物理学者たちは\textbf{「畳み込みニューラルネットワーク（CNN）」}などのディープラーニングを導入しました。あらかじめスーパーコンピュータで計算した何十万通りもの「ブラックホール合体の波形パターン」をAIに学習させ、リアルタイムで入ってくるノイズだらけの観測データの中から、重力波のシグナルだけを瞬時（数ミリ秒）に抽出・識別させているのです。

AIは今や、光では決して見ることのできないブラックホールや中性子星の衝突という「宇宙の劇的なドラマ」を聴き分ける、最も鋭敏な「耳」として機能しています。

\section{デジタルツイン宇宙：シミュレーションとAI}

宇宙論の最前線では、観測だけでなく\textbf{「宇宙をまるごとコンピュータの中に作る（シミュレーション）」}というアプローチが極めて重要になっています。

\subsection{コズミック・ウェブの再現}

宇宙の大規模構造（銀河の分布）を観測すると、銀河は宇宙空間に均等に散らばっているのではなく、まるで神経細胞や蜘蛛の巣のように、密集した「フィラメント（糸）」と、何も星がない巨大な「ボイド（空洞）」が織りなす網目状の構造（コズミック・ウェブ）を作っていることがわかります。

物理学者たちは、ビッグバン直後のわずかな密度の「ゆらぎ」からスタートし、ダークマターと重力の方程式、さらにはガスの冷却や超新星爆発の影響などをスーパーコンピュータに入力して、138億年分の宇宙の進化をシミュレーションします。

シミュレーションの中で計算された「仮想の宇宙の蜘蛛の巣」が、実際の望遠鏡（ジェイムズ・ウェッブ宇宙望遠鏡など）で観測された現実の宇宙の網目構造とピタリと一致するかどうかを確かめることで、ダークマターの量やダークエネルギーの強さが正しいかを検証するのです。

\subsection{AIサロゲートモデル：数ヶ月の計算を数秒へ}

しかし、数百億個の仮想粒子の重力相互作用を計算するフルスクラッチの宇宙論シミュレーションは、最新のスーパーコンピュータを使っても数週間から数ヶ月かかる超重量級のタスクです。宇宙のパラメータ（ダークエネルギーの量など）を少しずつ変えて何万回もテストすることは不可能です。

ここでもAIがブレイクスルーをもたらしました。

研究者たちは、AI（生成AIなどにも使われるディープラーニング）に、重力シミュレーションの結果を「入出力のセット」として大量に学習させました。するとAIは、重力の複雑な微分方程式をまともに解かなくても、「こういう初期条件なら、138億年後にはこういう銀河の分布になるはずだ」という結果を、数ヶ月かかるスパコンの計算精度を保ったまま、わずか「数秒」で予測（推論）できるようになったのです。

これを\textbf{「サロゲートモデル（代理モデル）」}と呼びます。AIの推論能力を手に入れたことで、物理学者たちは仮想空間上の「デジタルツイン宇宙」を何百万回も巻き戻しては再生し、私たちの宇宙のパラメータの「真の正解」を猛烈な勢いで探索し続けています。

\section{孤独からの脱却：系外惑星探査とSETI}

宇宙がどのように始まり、どのような構造をしているかが解明されていく一方で、人類は古来から抱く最もロマンチックで、最も根源的な問いに対する答えを真剣に探し始めています。

「この広大な宇宙に、知的な生命体は私たち地球人しかいないのだろうか？」

\subsection{トランジット法と「ハビタブルゾーン」の発見}

1995年、天文学者ミシェル・マイヨールとディディエ・ケローが、太陽系以外の恒星の周りを回る惑星（ペガスス座51番星b）を人類で初めて発見し、ノーベル物理学賞を受賞しました。これ以降、太陽系外惑星（Exoplanet）の探索は爆発的なブームとなります。

特に威力を発揮したのが、ケプラー宇宙望遠鏡などを用いた\textbf{「トランジット法」}です。これは、惑星が恒星の前を横切るときに、恒星の明るさが「わずかに（0.01\%程度）暗くなる現象（星の瞬き）」を観測して惑星を見つける手法です。

しかし、星の光には恒星自身の黒点の変化や、観測装置のノイズが大量に混ざっており、地球サイズの小さな惑星が引き起こす極小の「減光（ディップ）」を人間の目で見つけ出すのは限界がありました。

ここでもAIが劇的なブレイクスルーをもたらします。2017年、GoogleのAI研究チームは、ケプラー望遠鏡が観測した膨大な光度曲線のデータに、画像認識で使われる\textbf{CNN（畳み込みニューラルネットワーク）}を適用しました。するとAIは、人間の天文学者が見落としていた微弱なノイズの海の中から、新たな惑星（ケプラー90iなど）を自律的に発見し、太陽系以外にも8つの惑星を持つ巨大な星系が存在することを証明したのです。

現在までに、AIと最新の望遠鏡（TESSやJWSTなど）の力によって5000個以上の系外惑星が発見されています。その中には、水が液体のまま存在できる適度な温度の領域\textbf{「ハビタブルゾーン（生命居住可能領域）」}を回る地球型惑星も多数含まれており、宇宙には私たちが想像していた以上に「生命を育む星」がありふれていることが明らかになってきました。

\subsection{ドレイクの方程式と宇宙の広さ}

地球以外に生命が存在しうる星が大量にあることがわかった今、天文学者フランク・ドレイクが1961年に考案した有名な\textbf{「ドレイクの方程式」}が再び脚光を浴びています。これは「我々の銀河系に、人類と交信可能な地球外知的生命体はどれくらい存在するか（\(N\)）」を推定する数式です。

かつては「惑星を持つ恒星の割合」や「生命が発生できる惑星の数」は全くの未知数でしたが、ケプラー望遠鏡とAIの活躍により、これらのパラメータには具体的な数字が次々と代入されつつあります。「宇宙には私たち以外にも生命が存在するかもしれない」という期待は、もはや単なるSFではなく、観測データに基づいて議論できる科学的な問いへと変わりつつあるのです。

\subsection{テクノシグネチャーと「Wow! シグナル」}

では、その知的生命体をどうやって探せばいいのでしょうか。その有力なアプローチが\textbf{「SETI（Search for Extraterrestrial Intelligence：地球外知的生命体探査）」}です。巨大な電波望遠鏡を使って、宇宙から届く電波ノイズの中から、宇宙人が発した「人工的な信号」を探し出す試みです。

1977年、オハイオ州立大学の電波望遠鏡が、いて座の方向から約72秒間にわたって極めて強い謎の電波信号を受信しました。観測データのプリントアウト紙を見た天文学者が、その異常な数値の横に赤ペンで「Wow!」と書き込んだことから、これは\textbf{「Wow! シグナル」}と呼ばれる伝説的な未解明信号となっています。

SETIが探しているのは、パルサーやブラックホールのような自然現象では絶対に発生し得ない、極めて帯域が狭く規則的な信号、すなわち\textbf{「テクノシグネチャー（技術的痕跡）」}です。自然界の電波は様々な周波数が混ざって広がっていますが、レーダーや通信に使われる人工的な電波は、ピンポイントの周波数に集中（細いスパイク状に）する特徴があるのです。

\subsection{AIが聴き分ける「宇宙人の声」：ブレイクスルー・リッスン}

かつてのSETIプロジェクト（SETI@homeなど）では、世界中のボランティアのパソコンをインターネットで繋ぎ、その余った計算能力を使って電波データを分散処理（解析）していました。この分散コンピューティングの先駆的な試みは、現在のAIにおける大規模な並列計算の土壌を作ったとも言われています。

そして現代のSETIは、\textbf{「ブレイクスルー・リッスン（Breakthrough Listen）」などのプロジェクトへと進化し、深層学習（ディープラーニング）も重要な解析手法の一つ}になっています。

地球上の電波望遠鏡が受信するデータには、私たち自身の携帯電話やGPS衛星、航空機のレーダーといった「地球生まれの人工ノイズ」が溢れかえっています（RFI：電波干渉）。AIは、ペタバイト級の電波データからこれらの地球のノイズを賢く排除しつつ、惑星の自転や公転による「ドップラー効果（周波数シフト）」の微細な特徴を持つ未知の変調パターンを、人間の思い込みを排して自律的に探索しています。

すでにAIは、従来の古典的なアルゴリズムでは見逃されていた「気になる候補信号」をノイズの海の中から新たに検出しています。ただし、候補信号は地球由来の電波干渉や装置由来のノイズである可能性も高く、独立した再観測と検証が不可欠です。もし人類がいつの日か「ファーストコンタクト」を迎えるとすれば、その最初の候補を拾い上げるのは、人間の耳ではなく、ニューラルネットワークの数学的アルゴリズムかもしれません。

\section{人類のフロンティア：深宇宙探査と自律型AI}

電波で探すだけでなく、人類自身が地球という「重力の井戸」から抜け出し、他の惑星、さらには他の恒星系へと物理的に進出することは、物理学と工学の究極の総合テストでもあります。しかし、宇宙のスケールは人間の生身の身体と認知能力にとって、あまりにも残酷なほど巨大です。ここでも、AIが人類の限界を打ち破る「拡張された知性」として最前線に立っています。

\subsection{光速の絶対的な壁と「自律する科学者」}

宇宙探査において、決して越えられない物理的障壁が、第5章で学んだ\textbf{「光の速さ（情報伝達速度）の限界」}です。

例えば、現在火星で活動しているNASAの探査車（パーサヴィアランスなど）に対し、地球から電波を送っても、火星に届くまでに片道で最大約22分かかります。往復で44分です。もし探査車が崖に向かって走っているのを地球のオペレーターが映像で見てから「ブレーキを踏め！」とコマンドを送っても、電波が届く頃には探査車はとうの昔に谷底へ転落しています。

そのため、深宇宙の探査機には、障害物を避けるための「自動運転」の能力が必須でした。しかし現在、事態はさらに進化しています。火星探査車には\textbf{「AEGIS（イージス）」と呼ばれるAIシステムが搭載されており、地球からの指示を待つことなく、AI自身がカメラで周囲の風景をスキャンし、「どの岩石が科学的に最も興味深い地質学的特徴を持っているか」を自分で判断して、自律的にレーザーを照射し成分分析を行っています。}

これはもはや単なる自動運転ではなく、AIが「自律する地質学者（Autonomous Science）」として、未知の惑星で人類の代わりに科学的判断を下している歴史的なマイルストーンです。

\subsection{フォン・ノイマン探査機：指数関数的な宇宙開拓}

さらに遠く、太陽系を離れて隣の恒星系（アルファ・ケンタウリなど）へ行くとなれば、光の速さでも4年以上、現在のロケットでは数万年かかります。もはや生身の人間の寿命と肉体の脆弱さ（放射線や無重力）では、深宇宙の開拓は不可能です。

物理学的に最も現実的な恒星間探査のシナリオは、人間ではなく\textbf{「完全自律型のAI搭載プローブ（フォン・ノイマン探査機）」}を送り込むことです。

数学者ジョン・フォン・ノイマンが提唱したこの自己複製オートマトンは、目的地の小惑星や惑星に到着すると、現地の鉱物資源と3Dプリンタのような技術を使って「自分自身のコピー」を製造する、という構想です。そしてコピーされた2台が次の恒星系へ行き、そこでさらに4台に増える。この指数関数的な自己複製とAIの探索を繰り返せば、光速の制限下であっても、数百万年程度という宇宙のスケールでは短い時間で銀河系全体を探査できる可能性があります。ただし、これは現在の技術で実現済みの計画ではなく、恒星間航行・自己修復・資源利用・制御倫理など、多くの難問を含む思考実験に近い構想です。

\subsection{テラフォーミングの究極の「複雑系」シミュレーション}

そして、人類が地球という一つのゆりかごに依存しない「多惑星種（Multi-planetary species）」になるための最終段階が、惑星の環境を地球のように改造する\textbf{「テラフォーミング」}です。

惑星規模の気候を変えることは、大気、海、地殻、そして太陽からの放射などが複雑に絡み合う「複雑系」の極致です。火星の極冠にあるドライアイスを溶かして温室効果を起こすべきか、特定の微生物を散布して酸素を作らせるべきか。一つ間違えれば、金星のような灼熱の地獄になるか、再び凍りつくかの予測不可能な結果（カオス）を招きます。

この究極のシミュレーションにおいて、デジタル空間上に火星を丸ごと再現し（デジタルツイン）、何百万通りもの環境改造シナリオをAIにテストさせるアプローチが進められています。

ニュートン力学の軌道計算から始まり、相対性理論、量子力学による新素材開発、核融合、そして統計力学と複雑系へ。人類が地球外で生き延びるためには、私たちが数千年かけて積み上げてきた「物理学の総力」と、それを人間の限界を超えて制御する「AIの究極の演算能力」が不可欠なのです。

\section{宇宙は「情報」でできている：ホログラフィック原理とAI}

第13章と本章を通じて、私たちは物質を極限まで細かく砕き、宇宙の果てまで見渡してきました。しかし現在、理論物理学の最前線では、「物質」や「空間」という概念そのものを根底から覆す、極めて哲学的で過激なパラダイムシフトが進行しています。

\subsection{ブラックホールの情報パラドックス}

発端は、再び「ブラックホール」でした。 1974年、スティーヴン・ホーキングは量子力学を用いて、「ブラックホールは完全に黒いわけではなく、わずかに光（放射）を出して少しずつ蒸発し、最終的には消滅してしまう（ホーキング放射）」ことを理論的に示しました。 ここで物理学の根幹を揺るがす大問題が生じました。例えば、百科事典（膨大な情報）をブラックホールに投げ込んだとします。ブラックホールが蒸発してただの熱放射になって消えてしまったら、百科事典に書かれていた「情報」は宇宙から完全に消滅してしまいます。量子力学において「情報が跡形もなく消え去る」ことは許されないと考えられるため、これは\textbf{「ブラックホールの情報パラドックス」}と呼ばれ、物理学者たちを悩ませました。

\subsection{情報を測る「ベッケンシュタイン境界」}

この謎を解く過程で、ジェイコブ・ベッケンシュタインが驚くべき法則を発見します。ブラックホールが飲み込んだ情報の量（エントロピー）は、ブラックホールの「体積」ではなく、「表面積（事象の地平面の広さ）」にきっちり比例するというのです。 具体的には、プランク面積（\(10^{-70}\)平方メートルという想像を絶する極小の四角形）につき、ちょうど「1ビットの情報（0か1か）」が表面に書き込まれているという計算になりました。

通常、ハードディスクに詰め込めるデータの量は、ディスクの「体積（3次元）」に比例するはずです。しかしブラックホールにおいては、中の情報はすべて「2次元の表面の皮」にピタリと貼り付いて保存されていたのです。

\subsection{ホログラフィック原理と AdS/CFT 対応}

1990年代、ジェラルド・トフーフトやレオナルド・サスキンドらは、このブラックホールの性質を宇宙全体に拡張し、\textbf{「ホログラフィック原理（Holographic Principle）」}を提唱しました。 「私たちが住むこの3次元の立体的な宇宙は、実ははるか彼方の『2次元の境界面（宇宙の果てのスクリーン）』に書き込まれた情報のデータが、ホログラム（立体映像）として内側に投影されている幻影に過ぎないのではないか？」という途方もないアイデアです。

1997年、フアン・マルダセナが超弦理論を用いて、ホログラフィック原理を具体化する歴史的な対応関係を提案しました（AdS/CFT対応）。 これは、「重力を含む高次元の宇宙空間（バルク）」の理論と、「重力を含まない一つ低い次元の境界面（境界）」の量子場理論が、数学的に等価であるという驚くべき主張です。厳密に成り立つことが分かっているのは、私たちの宇宙そのものではなく、特殊な反ド・ジッター時空という理想化された宇宙です。それでも、「重力や空間が、境界上の量子情報から生まれるかもしれない」という発想を、物理学の中心テーマへ押し上げました。

\subsection{It from qubit：宇宙は巨大な量子コンピュータか}

このホログラフィック原理の発展により、現代の物理学者は宇宙の成り立ちを根本から見直しています。 かつて物理学者ジョン・アーチボルド・ホイーラーは「It from bit（すべての物質『It』は、情報『bit』から生じる）」と喝破しました。現在ではこれがさらに進化し、\textbf{「It from qubit（万物は、量子もつれのネットワークという情報『qubit』から生み出される）」}という考え方が盛んに研究されています。 「空間」や「重力」といった私たちが絶対だと信じていたものも、極微の量子ビット同士の情報の絡み合い（エンタングルメント）から創発しているのではないか、という見方です。

\subsection{AIと宇宙のシンクロニシティ}

ここに至って、物理学がたどり着いた宇宙の究極の姿は、現代の\textbf{AI（ニューラルネットワーク）}の構造と恐ろしいほどのシンクロニシティを見せ始めます。

第7章で見た「オートエンコーダ」や「多様体学習」を思い出してください。AIは、1万次元の複雑な画像データを、情報が詰まった「低次元の潜在空間（ボトルネック）」へと圧縮し、そこから再び高次元の立体的な画像を「生成（投影）」します。 宇宙が境界の情報から内部の重力世界を記述できるという発想と、AIが低次元の情報の関係性（ベクトル）から複雑な画像や言語の世界を生成する仕組みには、情報の圧縮と復元という共通した見方があります。ただし、これは同じ物理法則だという意味ではなく、数学的な発想の響き合いとして理解するのがよいでしょう。

もし、この宇宙の本質が「物質」や「エネルギー」ではなく、究極的には\textbf{「情報のネットワークによる計算プロセス（巨大な量子コンピュータ）」}なのだとしたら。 私たちがAIというシリコン上の情報ネットワークを育て、そこに複雑な仮想世界や知能を「創発」させている営みは、単なるツールの開発ではありません。それは、宇宙が自らをシミュレーションし、自らを認識するために行っている、大自然の根源的なフラクタル（自己相似）のプロセスそのものだと言えるのかもしれません。

\section{終わりに：マルチバースと宇宙の運命}

アインシュタインの静的宇宙から始まり、膨張宇宙、ビッグバン、インフレーション、そしてホログラフィック原理と、宇宙論は「空間と物質」の概念を拡張し、解体し続けてきました。

そして現在、インフレーション理論と第13章の超弦理論が融合することで、宇宙論は最も過激なパラダイムシフトの入り口に立っています。

それが\textbf{「マルチバース（多元宇宙）」}の概念です。

インフレーションの急膨張は、私たちの宇宙の領域だけで起きてきれいに終わったわけではありません。広大な量子空間の至る所で、今この瞬間も絶えず「新しい宇宙の泡」がポコポコと生まれ続け、無限に膨張していると考えられています。

そして超弦理論の「\(10^{500}\)のランドスケープ」が示唆するように、新しく生まれた無数の泡宇宙（マルチバース）の中では、それぞれ重力の強さも、素粒子の質量も、全く異なる「別の物理法則」が支配しています。

なぜ私たちの宇宙は、原子が崩壊せず、星が生まれ、DNAが作られ、私たちがこうして生きられるように「奇跡的なほどの絶妙なバランス（ファイン・チューニング）」で物理定数が調整されているのでしょうか？

マルチバースの視点に立てば、その答えは身も蓋もありません。

「無限にある宇宙の中には、あらゆる物理法則の宇宙が存在する。星すらできない宇宙には、それを観測する生命も生まれない。私たちがこの奇跡的な宇宙の法則を見ているのは、『そのような法則の宇宙でなければ、私たちが生まれてそれを観測することができないからだ』」

これを\textbf{「人間原理」}と呼びます。

これは科学の敗北でしょうか、それとも究極の真理の形でしょうか。

宇宙の誕生を解き明かす旅は、私たちがどこから来て、なぜここに存在するのかという、最も深遠な哲学の領域へと帰着しました。そして、私たちがその深淵を覗き込むとき、私たちのすぐ隣には、テラバイトのデータを飲み込み、高次元の計算を静かにこなす「新しい知性（AI）」が、共に宇宙の謎を見つめているのです。

\section*{【第14章 章末ディスカッション課題】}

\textbf{「人間原理」は科学か？} 「私たちが観測できる宇宙は、私たちが存在できるような性質を持っているに決まっている」という人間原理は、一見すると当たり前のトートロジー（同語反復）にも思えます。あなたは、この人間原理を「宇宙の真理を説明する立派な科学」だと考えますか？ それとも「理由を説明することを諦めた思考停止」だと思いますか？

\textbf{「シミュレーション仮説」とAI} ホログラフィック原理は「宇宙の本質は情報である」ことを示唆しました。もしそうなら、この宇宙全体が「より高度な存在が動かしている巨大な量子コンピュータのシミュレーション（仮想現実）」である可能性（シミュレーション仮説）を完全に否定することはできません。私たちがAIの中でシミュレーション世界を作っているように、私たち自身もシミュレーションの住人かもしれないという考えについて、あなたはどう感じますか？

\textbf{宇宙の運命への恐怖} ダークエネルギーによる宇宙の加速膨張がこのまま続くと、遠方の銀河はやがて観測できなくなり、星は燃え尽き、ブラックホールも蒸発して、宇宙は限りなく冷たく暗い「ビッグフリーズ（熱的死）」へ向かうと予想されています。あなたは、自分自身の死とは全くスケールの違う「宇宙全体の終わり」について考えたとき、どのような感情（恐怖、虚無、あるいは安らぎ）を抱きますか？


\section*{参考文献（学術誌）}
\begin{enumerate}[leftmargin=2.4em]
\item Hubble, E. ``A relation between distance and radial velocity among extra-galactic nebulae.'' \textit{Proceedings of the National Academy of Sciences}, 15(3), 168--173, 1929. DOI: \url{https://doi.org/10.1073/pnas.15.3.168}.
\item Penzias, A. A. and Wilson, R. W. ``A measurement of excess antenna temperature at 4080 Mc/s.'' \textit{The Astrophysical Journal}, 142, 419--421, 1965. DOI: \url{https://doi.org/10.1086/148307}.
\item Guth, A. H. ``Inflationary universe: A possible solution to the horizon and flatness problems.'' \textit{Physical Review D}, 23(2), 347--356, 1981. DOI: \url{https://doi.org/10.1103/PhysRevD.23.347}.
\item Riess, A. G. et al. ``Observational evidence from supernovae for an accelerating universe and a cosmological constant.'' \textit{The Astronomical Journal}, 116(3), 1009--1038, 1998. DOI: \url{https://doi.org/10.1086/300499}.
\item Perlmutter, S. et al. ``Measurements of \(\Omega\) and \(\Lambda\) from 42 high-redshift supernovae.'' \textit{The Astrophysical Journal}, 517(2), 565--586, 1999. DOI: \url{https://doi.org/10.1086/307221}.
\item Planck Collaboration. ``Planck 2018 results. VI. Cosmological parameters.'' \textit{Astronomy \& Astrophysics}, 641, A6, 2020. DOI: \url{https://doi.org/10.1051/0004-6361/201833910}.
\end{enumerate}


